\documentclass[11pt,a4paper]{article}

\usepackage[utf8]{inputenc}
\usepackage[T1]{fontenc}
\usepackage{lmodern}
\usepackage[french]{babel}
\usepackage[margin=2.3cm]{geometry}
\usepackage{xcolor}
\usepackage{booktabs}
\usepackage{array}
\usepackage{tikz}
\usetikzlibrary{arrows.meta,positioning,calc,shapes.geometric,fit,babel}
\usepackage{pgfplots}
\pgfplotsset{compat=1.16,/pgf/number format/.cd,use comma,1000 sep={\,}}
\usepackage{float}
\usepackage[most]{tcolorbox}
\usepackage{listings}
\usepackage{amsmath,amssymb}
\usepackage{caption}
\usepackage[hidelinks]{hyperref}

% Couleurs reprises de l'illustration interactive (atelier.html)
\definecolor{jzero}{HTML}{2F6FD6}
\definecolor{jun}{HTML}{D98A14}
\definecolor{jdeux}{HTML}{1F9D74}
\definecolor{accent}{HTML}{D1452F}
\definecolor{encre}{HTML}{18202A}
\definecolor{gris}{HTML}{5B6774}
\definecolor{grille}{HTML}{E6EAEF}
\definecolor{fond}{HTML}{EEF1F4}

\newtcolorbox{idee}[1][]{colback=fond,colframe=jzero,boxrule=0.6pt,arc=2pt,
  left=6pt,right=6pt,top=4pt,bottom=4pt,fonttitle=\bfseries,title=#1}
\newtcolorbox{calcul}{colback=fond,colframe=gris,boxrule=0.4pt,arc=2pt,
  left=6pt,right=6pt,top=3pt,bottom=3pt,fontupper=\ttfamily\small}
\newtcolorbox{pseudo}[1]{colback=white,colframe=encre,boxrule=0.5pt,arc=2pt,
  fonttitle=\bfseries\small,title=#1,left=4pt,right=4pt}

\lstdefinestyle{algo}{basicstyle=\small\ttfamily,columns=fullflexible,
  keepspaces=true,mathescape=true,numbers=left,numberstyle=\tiny\color{gris},
  numbersep=6pt,xleftmargin=14pt,
  keywordstyle=\bfseries\color{jzero},
  morekeywords={pour,tant,que,si,sinon,alors,fin,retourner,faire,chaque,dans,et,ou},
  literate={é}{{\'e}}1 {è}{{\`e}}1 {à}{{\`a}}1 {ê}{{\^e}}1 {ù}{{\`u}}1
           {ô}{{\^o}}1 {î}{{\^i}}1 {ç}{{\c c}}1 {É}{{\'E}}1 {û}{{\^u}}1
           {←}{{$\leftarrow$}}1 {≤}{{$\le$}}1 {×}{{$\times$}}1}
\lstdefinestyle{shell}{basicstyle=\small\ttfamily,columns=fullflexible,
  backgroundcolor=\color{fond},frame=none,xleftmargin=6pt,xrightmargin=6pt,
  keepspaces=true}
\lstdefinestyle{json}{basicstyle=\footnotesize\ttfamily,columns=fullflexible,
  backgroundcolor=\color{fond},frame=none,xleftmargin=6pt,keepspaces=true}
\lstdefinestyle{cpp}{language=C++,basicstyle=\footnotesize\ttfamily,
  columns=fullflexible,keepspaces=true,showstringspaces=false,
  keywordstyle=\bfseries\color{jzero},commentstyle=\itshape\color{jdeux!80!black},
  stringstyle=\color{jun!85!black},
  morekeywords={thread_local,nullptr,size_t,auto,constexpr},
  numbers=left,numberstyle=\tiny\color{gris},numbersep=7pt,
  xleftmargin=20pt,framexleftmargin=16pt,frame=l,rulecolor=\color{jzero!50},
  framerule=1.2pt,backgroundcolor=\color{fond},
  breaklines=true,breakatwhitespace=true,
  postbreak=\mbox{\textcolor{gris}{$\hookrightarrow$}\space},
  aboveskip=8pt,belowskip=6pt,
  captionpos=t,abovecaptionskip=3pt}
% Extrait reel du code source : \code{fichier}{premiere}{derniere}{titre}
% Les numeros affiches sont ceux du fichier (version du 19 octobre 2026).
\newcommand{\code}[4]{%
  \needspace{10\baselineskip}%
  \lstinputlisting[style=cpp,firstline=#2,lastline=#3,firstnumber=#2,
    title={\small\textbf{\texttt{\detokenize{#1}}}, lignes #2 à #3 --- #4}]{../#1}}
% Renvoi vers une ligne de l'extrait
\newcommand{\li}[1]{\textcolor{jzero}{\textbf{l.\,#1}}}
\newenvironment{lignes}{\begin{description}[style=multiline,leftmargin=2.1cm,font=\normalfont,itemsep=2pt]\small}{\end{description}}
\usepackage{enumitem}
\usepackage{needspace}

% ---------- Diagramme de Gantt pour l'atelier 3 x 3 ----------
% \op{machine}{debut}{fin}{couleur}{etiquette}
\newcommand{\op}[5]{%
  \fill[#4,rounded corners=1.5pt] (#2,{2-#1+0.13}) rectangle (#3,{2-#1+0.87});
  \node[white,font=\scriptsize\bfseries] at ({(#2+#3)/2},{2-#1+0.5}) {#5};}
% \ganttfond{T}
\newcommand{\ganttfond}[1]{%
  \foreach \t in {0,...,#1} \draw[grille] (\t,0) -- (\t,3);
  \foreach \m in {0,1,2} {
    \draw[gris!50] (0,{2-\m}) rectangle (#1,{3-\m});
    \node[anchor=east,font=\small\ttfamily,text=gris] at (-0.15,{2-\m+0.5}) {M\m};}
  \foreach \t in {0,2,...,#1} \node[font=\scriptsize\ttfamily,text=gris] at (\t,-0.3) {\t};}
\newcommand{\makespanligne}[1]{%
  \draw[accent,very thick,dashed] (#1,-0.05) -- (#1,3.25);
  \node[accent,font=\small\ttfamily\bfseries,fill=white,inner sep=1pt] at (#1,3.5) {makespan = #1};}

\title{\textbf{Résoudre le Job Shop par essaims particulaires}\\[4pt]
  \large Explication du code du projet \texttt{jsp\_pso}}
\author{Groupe B --- Optimisation combinatoire, HEH 2026--2027}
\date{}

\begin{document}
\maketitle

\begin{abstract}
Ce document explique, fichier par fichier, ce que fait le programme \texttt{jsp\_pso} :
comment il lit une instance, comment une \emph{particule} (une simple liste de nombres)
devient un planning, comment l'essaim améliore ces plannings au fil du temps, et
comment on l'empêche de stagner. Chaque mécanisme est illustré sur le petit atelier
$3\times3$ des slides du cours, puis mesuré sur des instances de référence (JSPLIB).
Le solveur est un PSO \emph{pur} : aucune recherche locale n'y est ajoutée.
\end{abstract}

\tableofcontents
\newpage

% =====================================================================
\section{Le problème : qui passe en premier sur la machine ?}

\subsection{L'atelier}

Un \emph{job shop} est un ensemble de \textbf{jobs} à exécuter sur un ensemble de
\textbf{machines}. Chaque job suit une \emph{recette} fixe : une suite d'opérations,
chacune sur une machine précise, pendant une durée précise. Trois contraintes doivent
être respectées :
\begin{enumerate}
  \item \textbf{ordre imposé} : les étapes d'un même job se font dans l'ordre de la recette ;
  \item \textbf{une machine, une opération à la fois} ;
  \item \textbf{pas d'interruption} : une étape commencée va jusqu'au bout.
\end{enumerate}
Les recettes ne sont pas négociables. \textbf{La seule décision, c'est l'ordre} dans
lequel chaque machine traite les jobs qui l'utilisent. L'objectif est de minimiser le
\textbf{makespan}, l'instant où le dernier job se termine.

\begin{table}[h]
\centering
\caption{L'exemple des slides : 3 jobs, 3 machines (machine $\cdot$ durée).}
\label{tab:recettes}
\begin{tabular}{lccc}
\toprule
 & Étape 1 & Étape 2 & Étape 3\\
\midrule
\textcolor{jzero}{\textbf{Job 0}} & M0 $\cdot$ 3 & M1 $\cdot$ 2 & M2 $\cdot$ 4\\
\textcolor{jun}{\textbf{Job 1}}   & M1 $\cdot$ 4 & M2 $\cdot$ 3 & M0 $\cdot$ 2\\
\textcolor{jdeux}{\textbf{Job 2}} & M2 $\cdot$ 2 & M0 $\cdot$ 5 & M1 $\cdot$ 3\\
\bottomrule
\end{tabular}
\end{table}

\subsection{Même atelier, deux ordres, deux résultats}

Plaçons les opérations une par une dans un ordre donné. Chaque opération démarre
\emph{au plus tôt} : quand son job est prêt \textbf{et} que sa machine est libre.

\begin{calcul}
début = max(job prêt, machine libre)\qquad fin = début + durée
\end{calcul}

L'ordre A ci-dessous est la suite $J_1, J_0, J_2, J_2, J_0, J_1, J_1, J_2, J_0$ : chaque
lettre désigne \og la prochaine étape de ce job \fg{}. Prenons la cinquième opération
placée, $J_0$ étape 2 sur M1 :

\begin{calcul}
J0 a fini son étape 1 à \ \ 3\\
M1 est occupée par J1 jusqu'à 4\\
début = max(3, 4) = 4 \quad $\Rightarrow$ J0 attend que la machine se libère\\
fin \ \ = 4 + 2 = 6
\end{calcul}

\begin{figure}[h]
\centering
\begin{tikzpicture}[x=0.66cm,y=0.72cm]
  \node[anchor=west,font=\bfseries] at (-1.3,4.3) {Ordre A (exemple 1 des slides)};
  \ganttfond{14}
  \op{0}{0}{3}{jzero}{J0·1} \op{0}{3}{8}{jdeux}{J2·2} \op{0}{8}{10}{jun}{J1·3}
  \op{1}{0}{4}{jun}{J1·1}   \op{1}{4}{6}{jzero}{J0·2} \op{1}{8}{11}{jdeux}{J2·3}
  \op{2}{0}{2}{jdeux}{J2·1} \op{2}{4}{7}{jun}{J1·2}   \op{2}{7}{11}{jzero}{J0·3}
  \makespanligne{11}
\end{tikzpicture}

\vspace{6mm}
\begin{tikzpicture}[x=0.66cm,y=0.72cm]
  \node[anchor=west,font=\bfseries] at (-1.3,4.3) {Ordre B (exemple 2 des slides)};
  \ganttfond{14}
  \op{0}{0}{3}{jzero}{J0·1} \op{0}{3}{8}{jdeux}{J2·2} \op{0}{12}{14}{jun}{J1·3}
  \op{1}{3}{5}{jzero}{J0·2} \op{1}{5}{9}{jun}{J1·1}   \op{1}{9}{12}{jdeux}{J2·3}
  \op{2}{0}{2}{jdeux}{J2·1} \op{2}{5}{9}{jzero}{J0·3} \op{2}{9}{12}{jun}{J1·2}
  \makespanligne{14}
\end{tikzpicture}
\caption{Les mêmes recettes avec deux ordres de placement. Dans l'ordre B, $J_0$ passe
en premier sur M1 : $J_1$ doit attendre jusqu'à 5 et tout le reste glisse
(\emph{effet domino}). Le makespan passe de 11 à 14.}
\label{fig:gantt-ab}
\end{figure}

\subsection{Pourquoi une métaheuristique ?}

Pour ce petit atelier, il existe $\frac{9!}{3!\,3!\,3!} = 1\,680$ ordres de placement.
Nous les avons tous énumérés : le meilleur donne \textbf{11}, le pire \textbf{28}, et
384 ordres atteignent l'optimum. Pour 10 jobs et 10 machines (l'exemple \texttt{10\_10}
des slides), le nombre d'ordres monte à $\frac{100!}{(10!)^{10}} \approx 2{,}4\times10^{92}$.
Il est impossible de tout essayer : il faut une méthode qui \emph{cherche intelligemment}.
C'est le rôle de l'optimisation par essaims particulaires.

% =====================================================================
\section{Vue d'ensemble du programme}

\begin{figure}[H]
\centering
\begin{tikzpicture}[
  box/.style={draw=encre,rounded corners=3pt,fill=white,align=center,font=\small,
              minimum height=1.05cm,inner xsep=5pt},
  file/.style={box,fill=fond},
  arr/.style={-{Stealth[length=2.2mm]},thick,encre},
  lab/.style={font=\scriptsize,text=gris},
  node distance=9mm and 12mm]
  \node[file] (in) {\texttt{10\_10}\\ \scriptsize fichier texte};
  \node[box,right=of in] (read) {\texttt{read\_instance}\\ \scriptsize instance.cpp};
  \node[box,right=of read,fill=jzero!10] (solve)
       {\texttt{solve}\\ \scriptsize pso.cpp : essaims};
  \node[box,right=of solve] (out) {\texttt{write\_schedule\_json}\\ \scriptsize schedule.cpp};
  \node[box,below=of solve,fill=jdeux!10] (dec)
       {\texttt{decode\_active}\\ \scriptsize schedule.cpp : clés $\to$ planning};
  \node[file,below=of out] (json) {\texttt{10\_10.json}\\ \scriptsize résultat};
  \node[box,above=of solve,fill=fond] (main) {\texttt{main}\\ \scriptsize main.cpp : options};
  \draw[arr] (in) -- (read);
  \draw[arr] (read) -- node[lab,below]{Instance} (solve);
  \draw[arr] ([xshift=-7mm]solve.south) -- node[lab,left]{particule} ([xshift=-7mm]dec.north);
  \draw[arr] ([xshift=7mm]dec.north) -- node[lab,right]{makespan} ([xshift=7mm]solve.south);
  \draw[arr] (solve) -- node[lab,below]{record} (out);
  \draw[arr] (out) -- (json);
  \draw[arr] (main) -- node[lab,right]{options} (solve);
\end{tikzpicture}
\caption{Le trajet des données dans le programme.}
\label{fig:archi}
\end{figure}

\begin{table}[h]
\centering
\caption{Les fichiers du projet.}
\begin{tabular}{>{\ttfamily}l p{10.5cm}}
\toprule
\textrm{Fichier} & Rôle\\
\midrule
src/instance.cpp & lit le fichier d'instance et vérifie qu'il est bien formé\\
src/schedule.cpp & les deux décodeurs (particule $\to$ planning) et l'écriture du JSON\\
src/pso.cpp      & l'essaim : mise à jour des particules, topologies, redémarrage, multi-essaims\\
src/main.cpp     & lit les options de la ligne de commande et lance le solveur\\
tests/test\_core.cpp & tests automatiques (lecture, exemple du cours, PSO)\\
tools/validate.py & vérificateur indépendant d'un JSON produit\\
\bottomrule
\end{tabular}
\end{table}

% =====================================================================
\section{Lire l'instance (\texttt{instance.cpp})}

Le fichier d'entrée a le format des slides : un nom, puis le nombre de jobs et de
machines, puis une ligne par job contenant les couples \texttt{machine durée}.

\begin{lstlisting}[style=shell]
10_10
10 10
4 88 8 68 6 94 5 99 1 67 2 89 9 77 7 99 0 86 3 92
...
\end{lstlisting}

En mémoire, l'atelier devient une structure \texttt{Instance} : le nom, \texttt{jobs},
\texttt{machines}, et \texttt{operations[job][étape]} qui contient le couple
\texttt{(machine, duration)}.

La lecture est volontairement \textbf{robuste}, car une erreur de lecture le jour du
concours coûterait une instance entière :
\begin{itemize}
  \item les lignes vides, les commentaires \texttt{\#} et les fins de ligne Windows
        (\texttt{\textbackslash r}) sont acceptés ;
  \item si la première ligne contient déjà les dimensions, le nom est pris dans le nom du
        fichier (le JSON de sortie doit porter le même nom que l'instance) ;
  \item un fichier mal formé est refusé avec un message clair : couple manquant ou en
        trop, machine inexistante, machine utilisée deux fois par un même job, job manquant ;
  \item des lignes en trop après le dernier job déclenchent seulement un avertissement.
\end{itemize}

\subsection{Dans le code}

\code{src/instance.hpp}{6}{16}{la structure qui représente l'atelier}
\begin{lignes}
  \item[\li{6--9}] une \texttt{Operation} = une étape de recette : sur quelle machine, pendant combien de temps.
  \item[\li{15}] \texttt{operations[job][étape]} : un tableau de recettes. Pour l'exemple $3\times3$,
        \texttt{operations[1][0]} vaut \texttt{\{machine 1, durée 4\}} (première étape du job 1).
\end{lignes}

\code{src/instance.cpp}{33}{61}{nettoyer le fichier et trouver le nom}
\begin{lignes}
  \item[\li{37--41}] on lit tout le fichier ligne par ligne, et on ne garde que les lignes utiles
        (la fonction \texttt{is\_comment\_or\_empty} écarte les lignes vides et celles qui
        commencent par \texttt{\#}).
  \item[\li{51--55}] on essaie de lire deux entiers sur la première ligne. Si elle ne contient
        \emph{que} deux entiers, c'est déjà la ligne des dimensions : le fichier n'a pas de nom.
  \item[\li{56}] dans ce cas, le nom est celui du fichier sans extension (\texttt{stem}) :
        \texttt{instances/abz5} donnera \texttt{abz5.json}.
  \item[\li{58--59}] sinon, la première ligne est le nom (\texttt{trim} retire les espaces et le
        \texttt{\textbackslash r} Windows), et on passe à la ligne suivante.
\end{lignes}

\code{src/instance.cpp}{75}{105}{lire chaque recette et la vérifier}
\begin{lignes}
  \item[\li{75}] une ligne par job.
  \item[\li{81--82}] \texttt{machine\_used} retient les machines déjà vues dans cette recette.
  \item[\li{85--90}] on lit un couple \texttt{machine durée}. Si la lecture échoue (pas assez de
        nombres, ou pas un nombre) ou si la machine n'existe pas, on arrête avec un message clair.
  \item[\li{92--96}] un job ne peut pas passer deux fois par la même machine.
  \item[\li{98}] l'étape est ajoutée à la recette du job.
  \item[\li{100--104}] s'il reste des nombres sur la ligne, elle est trop longue : erreur.
\end{lignes}

% =====================================================================
\section{Représenter une solution : les clés aléatoires}

Le PSO a été inventé pour des problèmes \emph{continus} : une particule y est un point
de $\mathbb{R}^D$ qui se déplace avec une vitesse. Or un ordonnancement est
\emph{discret}. Le pont entre les deux mondes est le codage par \textbf{clés aléatoires}.

\begin{idee}[Idée clé]
Une particule est une liste de $D = n \times m$ nombres réels : \textbf{une clé par
opération}. La clé de l'étape $o$ du job $j$ est rangée à l'indice $j\cdot m + o$.
Un \emph{décodeur} transforme ces nombres en planning, ce qui donne un makespan.
\end{idee}

Sur l'atelier $3\times3$, une particule contient donc 9 nombres. Avec le décodeur le plus
simple (\og semi-actif \fg{}, section~\ref{sec:semi}), on trie les clés par ordre
croissant ; chaque clé triée désigne un job, et on obtient une séquence de placement.

\begin{table}[h]
\centering
\caption{Une particule qui code l'ordre A. Trier les clés redonne la séquence
$J_1, J_0, J_2, J_2, J_0, J_1, J_1, J_2, J_0$, donc le planning de makespan 11.}
\label{tab:cles}
\small
\begin{tabular}{l*{9}{c}}
\toprule
Opération & J0·1 & J0·2 & J0·3 & J1·1 & J1·2 & J1·3 & J2·1 & J2·2 & J2·3\\
Indice $j\cdot m+o$ & 0 & 1 & 2 & 3 & 4 & 5 & 6 & 7 & 8\\
\midrule
Clé & 0,12 & 0,40 & 0,88 & 0,05 & 0,52 & 0,61 & 0,21 & 0,33 & 0,77\\
Rang après tri & 2 & 5 & 9 & 1 & 6 & 7 & 3 & 4 & 8\\
\bottomrule
\end{tabular}
\end{table}

L'intérêt est double : \emph{n'importe quelle} liste de nombres donne un planning valide
(les contraintes sont garanties par le décodeur), et l'on peut appliquer telles quelles
les formules continues du PSO.

\code{src/pso.cpp}{29}{45}{des clés à la séquence de jobs (décodeur \og semi \fg{})}
\begin{lignes}
  \item[\li{33--36}] \texttt{indices} = $0, 1, \dots, D-1$ : un numéro par opération.
  \item[\li{37--38}] on trie ces numéros \emph{selon la valeur de leur clé}. Sur la
        table~\ref{tab:cles}, on obtient $3, 0, 6, 7, 1, 4, 5, 8, 2$ (la clé 0,05 de l'indice 3
        est la plus petite, etc.). \texttt{stable\_sort} garde l'ordre des indices en cas
        d'égalité, pour que le résultat soit toujours le même.
  \item[\li{41--43}] chaque indice divisé par $m$ (division entière) donne son job :
        $3/3=1$, $0/3=0$, $6/3=2$\dots{} On retrouve la séquence
        $1, 0, 2, 2, 0, 1, 1, 2, 0$ de l'ordre A.
\end{lignes}

% =====================================================================
\section{Décoder une particule (\texttt{schedule.cpp})}

\subsection{Le décodeur semi-actif}
\label{sec:semi}

C'est le décodeur utilisé dans les exemples des slides : on prend la séquence de jobs
obtenue par le tri, et on place chaque opération \emph{au plus tôt} derrière celles déjà
placées. C'est exactement le calcul $\text{début}=\max(\text{job prêt},\text{machine libre})$
de la figure~\ref{fig:gantt-ab}.

\begin{pseudo}{decode\_semi\_active}
\begin{lstlisting}[style=algo]
job_prêt[j] ← 0 ; machine_libre[k] ← 0 ; suivante[j] ← 0
pour chaque job j dans la séquence faire
    (k, durée) ← operations[j][suivante[j]] ; suivante[j] ← suivante[j] + 1
    début ← max(job_prêt[j], machine_libre[k]) ; fin ← début + durée
    job_prêt[j] ← fin ; machine_libre[k] ← fin
retourner le plus grand fin   $\triangleright$ le makespan
\end{lstlisting}
\end{pseudo}

\code{src/schedule.cpp}{20}{38}{le décodeur semi-actif}
\begin{lignes}
  \item[\li{20--22}] trois compteurs : la prochaine étape de chaque job, la date où chaque job
        est prêt, la date où chaque machine se libère. Tout commence à 0.
  \item[\li{27--28}] \texttt{next[job]++} donne l'étape à placer \emph{puis} passe à la
        suivante ; \texttt{task} est cette étape (machine et durée).
  \item[\li{31}] c'est exactement l'encadré
        \texttt{début = max(job prêt, machine libre)} de la section 1.
  \item[\li{34--35}] le job et la machine sont tous deux occupés jusqu'à \texttt{end}.
  \item[\li{36}] le makespan est la plus grande date de fin rencontrée.
\end{lignes}

Son défaut : il ne remplit jamais un \og trou \fg{} laissé sur une machine. Beaucoup de
particules différentes donnent de mauvais plannings, et l'essaim perd son temps.

\subsection{Le décodeur Giffler-Thompson paramétré (par défaut)}

Le décodeur par défaut (\texttt{decode\_active}) suit l'algorithme de Giffler et Thompson,
avec un paramètre $\delta$ (Gonçalves \emph{et al.}). Il construit le planning
\emph{chronologiquement} et ne consulte la particule qu'au moment d'un \textbf{conflit}.
À chaque étape :

\begin{enumerate}
  \item pour chaque job, on calcule le début au plus tôt de sa prochaine opération et sa fin ;
  \item $\varphi$ = la plus petite fin possible. Sa machine $M^*$ est celle où se joue le
        prochain conflit ;
  \item $\sigma$ = le plus petit début possible sur $M^*$ ;
  \item les \textbf{candidats} sont les opérations sur $M^*$ qui peuvent commencer dans la
        fenêtre $[\sigma,\ \sigma + \delta(\varphi-\sigma)]$ ;
  \item \textbf{la particule décide} : le candidat dont la clé est la plus petite passe.
\end{enumerate}

\begin{figure}[h]
\centering
\begin{tikzpicture}[x=1.25cm,y=1cm]
  \draw[-{Stealth},gris] (-0.3,0) -- (8.4,0) node[right,font=\small] {temps};
  \foreach \t in {0,...,8} \draw[gris] (\t,0.08) -- (\t,-0.08) node[below,font=\scriptsize\ttfamily] {\t};
  \fill[jzero!18] (0,0.2) rectangle (0.6,1.9);
  \fill[accent!10] (0.6,0.2) rectangle (3,1.9);
  % les deux fenetres, en accolades au-dessus
  \draw[jzero,thick,|-|] (0,2.15) -- (0.6,2.15);
  \node[font=\scriptsize,text=jzero!80!black,anchor=east] at (-0.1,2.15) {fenêtre si $\delta=0{,}2$};
  \draw[accent,thick,|-|] (0,2.6) -- (3,2.6);
  \node[font=\scriptsize,text=accent,anchor=west] at (3.1,2.6) {fenêtre si $\delta=1$};
  \draw[very thick,jzero] (0,0.2) -- (0,1.9);
  \node[font=\scriptsize,anchor=east] at (-0.05,1.05) {$\sigma=0$};
  \draw[very thick,accent,dashed] (3,0.2) -- (3,1.9);
  \node[font=\scriptsize,text=accent,anchor=west,fill=white,inner sep=1pt] at (3.05,1.1) {$\varphi=3$};
  \fill[jzero,rounded corners=1.5pt] (0,1.3) rectangle (3,1.8);
  \node[white,font=\scriptsize\bfseries] at (1.5,1.55) {J0·1 : 0 $\to$ 3};
  \fill[jdeux,rounded corners=1.5pt] (2,0.35) rectangle (7,0.85);
  \node[white,font=\scriptsize\bfseries] at (4.5,0.6) {J2·2 : 2 $\to$ 7};
\end{tikzpicture}
\caption{Étape 2 du décodage de l'atelier $3\times3$ : J0 et J2 veulent tous deux la
machine M0. La fenêtre réglée par $\delta$ fixe qui est candidat. Avec $\delta=1$, les deux
sont candidats et c'est la clé de la particule qui tranche ; avec $\delta=0{,}2$, seul J0
peut commencer assez tôt (J2 ne commence qu'à $2 > 0{,}6$).}
\label{fig:fenetre}
\end{figure}

Le paramètre $\delta$ règle la taille de l'espace de recherche. Avec $\delta=1$, on
obtient des plannings \emph{actifs} : on ne peut avancer aucune opération sans en
retarder une autre, et l'optimum fait toujours partie de cet ensemble. Avec $\delta=0$,
on obtient des plannings \emph{sans retard} : une machine n'attend jamais quand une
opération est prête. Cet ensemble est plus petit et souvent très bon, mais il ne contient
pas toujours l'optimum. La valeur retenue, $\delta=0{,}2$, est un compromis mesuré
expérimentalement (section~\ref{sec:resultats}).

\begin{table}[h]
\centering
\caption{Déroulé complet du décodeur ($\delta=0{,}2$) sur l'atelier $3\times3$. Chaque
case de candidat donne la machine et l'intervalle début--fin au plus tôt.}
\label{tab:trace}
\footnotesize
\setlength{\tabcolsep}{4pt}
\begin{tabular}{c l l l c c c l}
\toprule
Pas & \textcolor{jzero}{J0} & \textcolor{jun}{J1} & \textcolor{jdeux}{J2} & $\varphi$ ($M^*$) & $\sigma$ & fenêtre & placé\\
\midrule
1 & M0 0--3  & M1 0--4  & M2 0--2  & 2 (M2)  & 0 & $[0;\,0{,}4]$ & J2 : 0--2\\
2 & M0 0--3  & M1 0--4  & M0 2--7  & 3 (M0)  & 0 & $[0;\,0{,}6]$ & J0 : 0--3\\
3 & M1 3--5  & M1 0--4  & M0 3--8  & 4 (M1)  & 0 & $[0;\,0{,}8]$ & J1 : 0--4\\
4 & M1 4--6  & M2 4--7  & M0 3--8  & 6 (M1)  & 4 & $[4;\,4{,}4]$ & J0 : 4--6\\
5 & M2 6--10 & M2 4--7  & M0 3--8  & 7 (M2)  & 4 & $[4;\,4{,}6]$ & J1 : 4--7\\
6 & M2 7--11 & M0 7--9  & M0 3--8  & 8 (M0)  & 3 & $[3;\,4]$     & J2 : 3--8\\
7 & M2 7--11 & M0 8--10 & M1 8--11 & 10 (M0) & 8 & $[8;\,8{,}4]$ & J1 : 8--10\\
8 & M2 7--11 & --       & M1 8--11 & 11 (M2) & 7 & $[7;\,7{,}8]$ & J0 : 7--11\\
9 & --       & --       & M1 8--11 & 11 (M1) & 8 & $[8;\,8{,}6]$ & J2 : 8--11\\
\bottomrule
\end{tabular}
\end{table}

Sur ce petit exemple, toutes les fenêtres ne contiennent qu'un seul candidat : le
décodeur seul construit directement le planning optimal (makespan 11), le même que
l'ordre A de la figure~\ref{fig:gantt-ab}. Sur une vraie instance $10\times10$, les
conflits à plusieurs candidats sont nombreux, et ce sont eux que l'essaim apprend à
trancher.

\subsubsection*{Dans le code}

La boucle principale de \texttt{decode\_active} fait une opération par tour
(\li{70} : $n\times m$ tours en tout). Chaque tour se lit en trois temps, qui
correspondent aux colonnes de la table~\ref{tab:trace}.

\code{src/schedule.cpp}{74}{89}{temps 1 : trouver $\varphi$, $M^*$ et les $\sigma$}
\begin{lignes}
  \item[\li{74--75}] on remet à \og l'infini \fg{} le plus petit début connu de chaque machine.
  \item[\li{78--80}] on parcourt les jobs ; \texttt{current\_machine[job]} vaut $-1$ quand le
        job est terminé, on le saute.
  \item[\li{81}] début au plus tôt de la prochaine étape du job : toujours le même
        \texttt{max(job prêt, machine libre)}.
  \item[\li{83}] on retient au passage le plus petit début sur chaque machine : ce sera
        $\sigma$ pour la machine du conflit.
  \item[\li{84--88}] on cherche la plus petite date de fin : c'est $\varphi$, et sa machine
        devient \texttt{conflict\_machine} ($M^*$). Au pas 1 du 3$\times$3, c'est J2 sur M2 qui
        finit le plus tôt (2), donc $\varphi=2$ et $M^*=$ M2.
\end{lignes}

\code{src/schedule.cpp}{90}{106}{temps 2 : la fenêtre et le choix de la particule}
\begin{lignes}
  \item[\li{92--93}] la fenêtre $[\sigma,\ \sigma+\delta(\varphi-\sigma)]$ : \texttt{limit}
        en est la borne droite. Au pas 2 : $\sigma=0$, $\varphi=3$, $\delta=0{,}2$, donc
        \texttt{limit} $=0{,}6$.
  \item[\li{97--100}] on écarte les jobs qui ne veulent pas $M^*$, ou qui ne peuvent pas
        commencer dans la fenêtre (au pas 2, J2 commence à $2 > 0{,}6$ : il est écarté).
  \item[\li{101}] la clé de la particule pour \emph{cette} opération : indice
        \texttt{job * machines + next[job]}, comme dans la table~\ref{tab:cles}.
  \item[\li{102--105}] parmi les candidats, la plus petite clé gagne. \textbf{C'est la seule
        ligne où la particule influence le planning.}
\end{lignes}

\code{src/schedule.cpp}{107}{115}{temps 3 : placer l'opération choisie}
\begin{lignes}
  \item[\li{107--109}] l'opération choisie commence à son début au plus tôt ; on calcule sa fin.
  \item[\li{110}] si l'appelant veut le planning complet (nouveau record), on note la date de début.
  \item[\li{111--112}] le job et la machine sont occupés jusqu'à cette fin.
  \item[\li{113--114}] le job passe à sa machine suivante, ou à $-1$ s'il a fini sa recette.
\end{lignes}

\paragraph{Vitesse.} Le décodeur est appelé des millions de fois. Ses tableaux de
travail sont réutilisés d'un appel à l'autre (\texttt{thread\_local}) et il ne fait que
deux passes sur les jobs par opération placée. Il ne construit le planning complet
(\texttt{schedule\_from\_starts}) que lorsqu'un nouveau record est trouvé. Résultat :
environ $170\,000$ évaluations par seconde sur un $10\times10$, et $46\,\mu$s par
évaluation sur un $30\times20$.

\code{src/schedule.cpp}{53}{59}{des tableaux de travail réutilisés}
\begin{lignes}
  \item[\li{53--59}] \texttt{thread\_local} : chaque thread garde ses propres tableaux d'un
        appel à l'autre. On évite ainsi de réserver de la mémoire des millions de fois, sans
        que deux essaims parallèles écrivent dans les mêmes tableaux.
\end{lignes}

\code{src/pso.cpp}{49}{63}{évaluer une particule}
\begin{lignes}
  \item[\li{51--57}] avec \texttt{-{}-decoder semi}, on trie les clés puis on applique le
        décodeur semi-actif.
  \item[\li{58}] cas le plus fréquent : on ne veut que le makespan, sans construire le planning.
  \item[\li{59--61}] pour un nouveau record seulement, on récupère les dates de début et on
        construit le planning complet qui ira dans le JSON.
\end{lignes}

% =====================================================================
\section{L'essaim particulaire (\texttt{pso.cpp})}

\subsection{Le principe}

L'optimisation par essaims particulaires (Kennedy et Eberhart, 1995) imite une nuée
d'oiseaux. Chaque particule $i$ possède :
\begin{itemize}
  \item une \textbf{position} $x_i$ (ses $D$ clés, donc un planning) ;
  \item une \textbf{vitesse} $v_i$ ;
  \item son \textbf{meilleur souvenir} $p_i$ (\emph{pbest}), la meilleure position qu'elle a visitée ;
  \item un \textbf{guide social} $g_i$ : le meilleur souvenir de ses voisins (\emph{gbest}
        quand le voisinage est tout l'essaim).
\end{itemize}

À chaque tour, chaque dimension $d$ de chaque particule est mise à jour :
\[
v_{i,d} \leftarrow \underbrace{w\,v_{i,d}}_{\text{inertie}}
 + \underbrace{c_1 r_1 \,(p_{i,d}-x_{i,d})}_{\text{mémoire}}
 + \underbrace{c_2 r_2 \,(g_{i,d}-x_{i,d})}_{\text{influence sociale}},
\qquad
x_{i,d} \leftarrow x_{i,d} + v_{i,d}
\]
où $r_1, r_2$ sont tirés au hasard dans $[0,1]$ à chaque fois.

\code{src/pso.cpp}{18}{27}{une particule en mémoire}
\begin{lignes}
  \item[\li{19--21}] $x_i$, $v_i$ et $p_i$ : trois listes de $D$ nombres.
  \item[\li{22}] le makespan de $p_i$, pour ne pas le recalculer.
  \item[\li{25--26}] utilisés seulement par la variante CLPSO (section~\ref{sec:clpso}).
\end{lignes}

\code{src/pso.cpp}{310}{324}{la formule du PSO, telle qu'elle est codée}
\begin{lignes}
  \item[\li{310--311}] \texttt{social} est $g_i$, le meilleur souvenir du guide de la
        particule (choisi par la topologie, section~\ref{sec:anneau}).
  \item[\li{314}] la mise à jour se fait dimension par dimension (\texttt{key} $= d$).
  \item[\li{315}] $w\,v_{i,d}$ : l'inertie.
  \item[\li{316--318}] $c_1 r_1 (p_{i,d}-x_{i,d})$ : la mémoire. \texttt{unit(random)} tire
        un nouveau $r_1$ dans $[0,1]$ pour chaque dimension.
  \item[\li{319--320}] $c_2 r_2 (g_{i,d}-x_{i,d})$ : l'influence sociale.
  \item[\li{321}] la vitesse est bornée à $\pm$\texttt{vmax}.
  \item[\li{322--323}] on mémorise la vitesse et on déplace la particule.
\end{lignes}

\begin{figure}[h]
\centering
\begin{tikzpicture}[scale=1.15,
  pt/.style={circle,fill=#1,inner sep=2.2pt},
  vec/.style={-{Stealth[length=2.6mm]},very thick,#1}]
  \node[pt=gris,label={left:{$x_i$ (position)}}] (x) at (0,0) {};
  \node[pt=jun,label={above:{$p_i$ (pbest)}}] (p) at (1.6,3.0) {};
  \draw[jun,dotted] (2.6,1.4) -- (p);
  \node[draw=encre,circle,very thick,minimum size=5.5mm,label={above:{$g_i$ (gbest)}}] (g) at (6.6,1.4) {};
  \fill[encre] (g) circle (1.3pt);
  \draw[jzero,dotted] (4.4,1.4) -- (g);
  \draw[vec=gris!80] (0,0) -- (2.0,-1.0) node[midway,below left,font=\small] {$w\,v_i$ (inertie)};
  \draw[vec=jun] (2.0,-1.0) -- (2.6,1.4) node[pos=0.3,right,font=\small] {mémoire};
  \draw[vec=jzero] (2.6,1.4) -- (4.4,1.4) node[midway,above,font=\small] {sociale};
  \draw[vec=accent,dashed] (0,0) -- (4.4,1.4) node[pos=0.28,above,sloped,font=\small] {nouvelle vitesse};
  \node[pt=accent] at (4.4,1.4) {};
  \node[font=\small,text=accent,anchor=north west] at (4.5,1.25) {nouvelle position};
\end{tikzpicture}
\caption{La nouvelle vitesse est la somme de trois tendances : continuer sur sa lancée,
revenir vers son meilleur souvenir, et aller vers le meilleur de ses voisins.}
\label{fig:vitesse}
\end{figure}

\subsection{Les réglages retenus}

\begin{itemize}
  \item \textbf{Coefficients} $c_1=c_2=1{,}49618$ : les valeurs de \og constriction \fg{} de
        Clerc et Kennedy, qui garantissent un essaim stable.
  \item \textbf{Inertie décroissante} de $w=0{,}9$ à $w=0{,}4$ au fil du temps : l'essaim
        explore au début et affine à la fin (\texttt{-{}-no-decreasing-inertia} fixe
        $w=0{,}7298$).
  \item \textbf{Vitesse bornée} à $\pm0{,}2$ par dimension (\texttt{-{}-vmax}) : une particule
        ne saute pas d'un bout à l'autre de l'espace en un tour.
  \item \textbf{Positions non bornées} : seul l'ordre relatif des clés compte. Si on
        bloquait les clés dans $[0,1]$, beaucoup s'empileraient à 0 ou à 1, deviendraient
        égales, et la particule perdrait de l'information.
  \item \textbf{Plateaux acceptés} : une position de \emph{même} makespan remplace le
        meilleur souvenir. Le JSP a de larges plateaux de solutions équivalentes ; sans
        cela, l'essaim s'y figerait.
  \item \textbf{Départ structuré} : dans le premier essaim, trois particules partent de
        règles de priorité classiques (SPT : opération la plus courte ; MWKR : job avec le
        plus de travail restant ; MOPNR : job avec le plus d'opérations restantes). Les
        autres partent au hasard.
\end{itemize}

\code{src/pso.cpp}{285}{290}{l'inertie décroissante}
\begin{lignes}
  \item[\li{287--288}] \texttt{ratio} va de 0 (début) à 1 (fin du budget).
  \item[\li{289--290}] $w = 0{,}9 - 0{,}5\times$\texttt{ratio} : 0,9 au départ, 0,4 à la fin.
        Avec \texttt{-{}-no-decreasing-inertia}, $w$ reste à 0,7298.
\end{lignes}

\code{src/pso.cpp}{179}{189}{tirer une particule au hasard}
\begin{lignes}
  \item[\li{181}] chaque clé est tirée uniformément dans $[0,1]$.
  \item[\li{182}] chaque vitesse est tirée dans $[-\texttt{vmax}, \texttt{vmax}]$.
  \item[\li{184--186}] au départ, le meilleur souvenir est la position elle-même.
  \item[\li{187--188}] remise à zéro de l'état CLPSO (utile après un redémarrage).
\end{lignes}

\code{src/pso.cpp}{72}{89}{les règles de priorité du départ}
\begin{lignes}
  \item[\li{73}] SPT : on préfère l'opération la plus courte (score $= -$durée, le plus grand score gagne).
  \item[\li{74--79}] MWKR additionne les durées restantes du job ; MOPNR compte les étapes restantes.
  \item[\li{81--89}] on construit une séquence en prenant à chaque pas le job au meilleur score.
        Les lignes 90 à 97 transforment cette séquence en clés (le rang divisé par $D$).
\end{lignes}

\code{src/pso.cpp}{226}{240}{créer l'essaim}
\begin{lignes}
  \item[\li{230}] toutes les particules sont d'abord tirées au hasard.
  \item[\li{233--238}] dans le premier essaim seulement, les particules 0, 1 et 2 sont
        remplacées par SPT, MWKR et MOPNR. Les autres essaims restent entièrement aléatoires,
        pour être plus divers.
  \item[\li{239}] chaque particule propose son planning comme record global.
\end{lignes}

\subsection{La boucle principale}

\begin{pseudo}{Un essaim (un thread)}
\begin{lstlisting}[style=algo]
initialiser 30 particules (3 règles de priorité + hasard) ; évaluer chacune
tant que le temps n'est pas écoulé faire
    w ← 0,9 - 0,5 × (temps écoulé / budget)
    calculer le guide de chaque particule (voisinage en anneau ou CLPSO)
    pour chaque particule i faire
        mettre à jour v$_i$ et x$_i$, dimension par dimension
        f ← makespan(décoder(x$_i$))
        si f ≤ f(p$_i$) alors p$_i$ ← x$_i$ ; publier si c'est un record global
    si pas de progrès depuis 200 tours alors relancer l'essaim (sauf la meilleure)
    tous les 50 tours : la pire particule adopte le record global (migration)
\end{lstlisting}
\end{pseudo}

Les lignes 282 à 376 de \texttt{pso.cpp} codent cette boucle. La mise à jour des vitesses
(lignes 7--8 du pseudo-code) a été montrée plus haut ; voici l'évaluation et la mémoire
(lignes 9--10).

\code{src/pso.cpp}{326}{339}{évaluer la nouvelle position et mettre à jour pbest}
\begin{lignes}
  \item[\li{326--327}] la nouvelle position est décodée en planning ; on récupère son makespan.
  \item[\li{328--332}] \texttt{stale} compte les tours sans amélioration stricte (pour CLPSO).
  \item[\li{335}] \texttt{<=} et non \texttt{<} : à makespan égal, la particule accepte quand
        même sa nouvelle position comme meilleur souvenir. Elle peut ainsi avancer sur un plateau.
  \item[\li{338}] si c'est mieux que le record de tous les essaims, \texttt{publish} l'enregistre
        (section~\ref{sec:parallele}).
\end{lignes}

% =====================================================================
\section{Empêcher l'essaim de stagner}

Un PSO classique converge \textbf{trop vite} : toutes les particules s'agglutinent autour
du premier bon planning trouvé, leurs vitesses tombent à zéro et plus rien ne bouge. Sur
abz5, la version de départ ne progressait plus après une seconde, sur les 55 disponibles.
Quatre mécanismes combattent ce phénomène. Ils restent tous des mécanismes d'essaim.

\begin{figure}[h]
\centering
\begin{tikzpicture}[
  p/.style={circle,draw=encre,fill=white,inner sep=0pt,minimum size=4.2mm,font=\tiny},
  best/.style={p,fill=accent!25,very thick}]
  % Topologie globale
  \begin{scope}
    \foreach \i in {0,...,7} \node[p] (g\i) at ({90+45*\i}:1.15) {\i};
    \foreach \i in {0,...,7} \foreach \j in {0,...,7} {
      \ifnum\i<\j \draw[gris!45] (g\i) -- (g\j); \fi}
    \node[font=\small\bfseries] at (0,-1.75) {globale};
    \node[font=\scriptsize,text=gris,align=center] at (0,-2.2) {tout le monde suit\\ le meilleur};
  \end{scope}
  % Anneau
  \begin{scope}[xshift=4.2cm]
    \foreach \i in {0,...,7} \node[p] (r\i) at ({90+45*\i}:1.15) {\i};
    \foreach \i in {0,...,7} {
      \pgfmathtruncatemacro{\a}{mod(\i+1,8)}
      \pgfmathtruncatemacro{\b}{mod(\i+2,8)}
      \draw[jzero,thick] (r\i) -- (r\a);
      \draw[jzero!40] (r\i) to[bend left=12] (r\b);}
    \node[font=\small\bfseries] at (0,-1.75) {anneau ($k=2$)};
    \node[font=\scriptsize,text=gris,align=center] at (0,-2.2) {chacun ne voit que\\ 2 voisins de chaque côté};
  \end{scope}
  % CLPSO
  \begin{scope}[xshift=8.6cm]
    \node[best] (c) at (0,0) {$x_i$};
    \foreach \i/\ang/\lab in {1/90/$d_1$,2/20/$d_2$,3/-50/$d_3$,4/-130/$d_4$,5/160/$d_5$} {
      \node[p] (q\i) at (\ang:1.25) {};
      \draw[-{Stealth},jdeux,thick] (q\i) -- (c) node[pos=0.35,font=\tiny,fill=white,inner sep=0.5pt] {\lab};}
    \node[font=\small\bfseries] at (0,-1.75) {CLPSO};
    \node[font=\scriptsize,text=gris,align=center] at (0,-2.2) {chaque dimension apprend\\ d'une particule différente};
  \end{scope}
\end{tikzpicture}
\caption{Les trois façons de choisir le guide social d'une particule.}
\label{fig:topologies}
\end{figure}

\subsection{La topologie en anneau}\label{sec:anneau}

Les particules sont disposées en cercle. Le guide social de chacune est le meilleur
souvenir parmi elle-même et ses $k=2$ voisins de chaque côté (\texttt{-{}-neighbors}). Un
bon planning se propage donc de proche en proche au lieu de contaminer tout l'essaim
d'un coup : plusieurs zones de l'espace restent explorées en parallèle plus longtemps.

\code{src/pso.cpp}{245}{271}{calculer le guide de chaque particule}
\begin{lignes}
  \item[\li{246--251}] on repère la meilleure particule de l'essaim (\texttt{best}).
  \item[\li{253--256}] topologie globale : tout le monde a \texttt{best} pour guide.
  \item[\li{259--267}] anneau : on regarde les particules de \texttt{index - k} à
        \texttt{index + k} et on garde celle qui a le meilleur souvenir.
  \item[\li{262}] le modulo referme le cercle : avec 30 particules, les voisines de la
        particule 0 sont 28, 29, 1 et 2.
\end{lignes}

\subsection{Le redémarrage partiel}

Si le meilleur de l'essaim ne s'améliore plus pendant 200 tours (\texttt{-{}-restart}),
l'essaim a convergé. On garde sa meilleure particule et on relance toutes les autres au
hasard, avec de nouvelles vitesses. Elles repartent explorer d'autres zones, et le
meilleur planning n'est jamais perdu.

\code{src/pso.cpp}{342}{348}{mesurer la stagnation}
\begin{lignes}
  \item[\li{343--345}] le meilleur de l'essaim a progressé : le compteur repart à zéro.
  \item[\li{347}] sinon, un tour de plus sans progrès.
\end{lignes}

\code{src/pso.cpp}{365}{375}{le redémarrage partiel}
\begin{lignes}
  \item[\li{366--367}] seulement pour les essaims en anneau, après \texttt{restart} = 200 tours
        sans progrès.
  \item[\li{370}] \texttt{keep} est la meilleure particule de l'essaim.
  \item[\li{371--373}] toutes les autres sont retirées au hasard avec \texttt{randomize}.
\end{lignes}

Nous avons aussi essayé de relancer les particules \emph{autour} de la meilleure
(position du record plus un bruit gaussien). C'était moins bon : l'essaim a besoin de vraie
diversité, pas d'une variation sur le même thème.

\subsection{CLPSO : l'apprentissage global}\label{sec:clpso}

Dans le PSO à apprentissage global (\emph{Comprehensive Learning PSO}, Liang \emph{et al.},
2006), une particule n'a plus de guide unique. Pour \textbf{chaque dimension} $d$, elle
choisit un \og exemple \fg{} : avec une probabilité $P_c(i)$, le meilleur souvenir d'une
autre particule tirée par tournoi (la meilleure de deux tirées au hasard), sinon le sien.
\[
v_{i,d} \leftarrow w\,v_{i,d} + c\, r\,\bigl(p_{\,e_i(d),\,d} - x_{i,d}\bigr),
\qquad
P_c(i) = 0{,}05 + 0{,}45\,\frac{e^{10\,i/(N-1)}-1}{e^{10}-1},
\qquad c = 1{,}49445
\]
Les exemples sont tirés à nouveau quand la particule ne progresse plus pendant 7 tours.
Comme chaque particule mélange des morceaux de plusieurs autres, l'essaim ne converge
jamais d'un bloc. C'est plus lent, mais cela explore beaucoup mieux. CLPSO entretient
lui-même sa diversité : on ne le redémarre pas.

\code{src/pso.cpp}{195}{224}{choisir les exemples d'une particule}
\begin{lignes}
  \item[\li{197--199}] la probabilité d'apprentissage $P_c(i)$ : 0,05 pour la première
        particule, 0,5 pour la dernière. Certaines particules restent fidèles à elles-mêmes,
        d'autres empruntent beaucoup.
  \item[\li{201--210}] le tournoi : deux autres particules tirées au hasard, on garde celle
        qui a le meilleur souvenir.
  \item[\li{211}] par défaut, chaque dimension apprend de la particule elle-même.
  \item[\li{213--218}] pour chaque dimension, avec la probabilité $P_c(i)$, l'exemple devient
        le gagnant d'un tournoi.
  \item[\li{219--222}] si aucune dimension n'a été tirée, on en force une : la particule doit
        apprendre de quelqu'un.
\end{lignes}

\code{src/pso.cpp}{295}{308}{la mise à jour CLPSO}
\begin{lignes}
  \item[\li{296--298}] on retire les exemples quand la particule stagne depuis 7 tours.
  \item[\li{300--301}] la cible de la dimension \texttt{key} est le meilleur souvenir de
        \emph{son} exemple : des dimensions différentes suivent des particules différentes.
  \item[\li{302--304}] un seul terme d'attraction au lieu de deux, avec $c=1{,}49445$.
\end{lignes}

\subsection{Le multi-essaims hétérogène}

Le programme lance \textbf{un essaim par cœur} du processeur (\texttt{-{}-threads 0}). Avec
la topologie \texttt{mixed}, un essaim sur deux tourne en anneau avec redémarrage, les
autres en CLPSO. Les mesures de la section~\ref{sec:resultats} montrent que CLPSO est
meilleur jusqu'à $15\times15$ et l'anneau au-delà. Le mélange est donc robuste quelle que
soit la taille des instances du concours.

\begin{figure}[h]
\centering
\begin{tikzpicture}[
  sw/.style={draw=encre,rounded corners=3pt,minimum width=2.3cm,minimum height=1cm,
             align=center,font=\small},
  arr/.style={-{Stealth[length=2mm]},thick}]
  \node[draw=accent,very thick,rounded corners=4pt,fill=accent!8,minimum width=3.2cm,
        minimum height=1.2cm,align=center] (rec) at (0,0)
        {\textbf{record global}\\ \scriptsize planning + position};
  \node[sw,fill=jzero!12] (s0) at (-4.2,1.5) {essaim 0\\ \scriptsize anneau};
  \node[sw,fill=jdeux!12] (s1) at (4.2,1.5)  {essaim 1\\ \scriptsize CLPSO};
  \node[sw,fill=jzero!12] (s2) at (-4.2,-1.5) {essaim 2\\ \scriptsize anneau};
  \node[sw,fill=jdeux!12] (s3) at (4.2,-1.5)  {essaim 3\\ \scriptsize CLPSO};
  \foreach \s in {s0,s1,s2,s3} {
    \draw[arr,encre] (\s) to[bend left=22] node[pos=0.4,font=\tiny,fill=white,inner sep=1pt]{publier} (rec);
    \draw[arr,accent,dashed] (rec) to[bend left=22] node[pos=0.6,font=\tiny,fill=white,inner sep=1pt]{migrer} (\s);}
  \node[font=\scriptsize,text=gris] at (0,-2.6) {\dots{} jusqu'à un essaim par cœur (12 sur la machine de test)};
\end{tikzpicture}
\caption{Les essaims tournent en parallèle et ne partagent que le record. Toutes les 50
itérations (\texttt{-{}-migrate}), la pire particule de chaque essaim adopte la position
du record global.}
\label{fig:multi}
\end{figure}

\code{src/pso.cpp}{277}{281}{quelle variante pour quel essaim ?}
\begin{lignes}
  \item[\li{280--281}] en mode \texttt{mixed}, les essaims de numéro impair (1, 3, 5\dots)
        tournent en CLPSO, les essaims pairs en anneau.
\end{lignes}

\code{src/pso.cpp}{350}{363}{la migration}
\begin{lignes}
  \item[\li{350--351}] tous les 50 tours, et seulement si un autre essaim a fait mieux que celui-ci.
  \item[\li{353--358}] on cherche la pire particule de l'essaim.
  \item[\li{359}] on verrouille le record global pendant qu'on le copie (un autre thread
        pourrait le modifier au même moment).
  \item[\li{360--362}] la pire particule prend la position et le makespan du record : elle va
        attirer ses voisines vers cette zone.
\end{lignes}

\code{src/pso.cpp}{378}{388}{lancer les essaims en parallèle}
\begin{lignes}
  \item[\li{380--381}] \texttt{-{}-threads 0} : autant d'essaims que de cœurs
        (\texttt{hardware\_concurrency}).
  \item[\li{383--385}] chaque thread exécute la fonction \texttt{worker}, c'est-à-dire un
        essaim complet, avec son numéro.
  \item[\li{386--388}] on attend que tous les essaims aient fini (fin du budget).
\end{lignes}

% =====================================================================
\section{Parallélisme et fiabilité}\label{sec:parallele}

\begin{itemize}
  \item \textbf{Partage du record.} Le makespan du record est un \texttt{std::atomic}, lu
        sans verrou à chaque évaluation. Le planning et la position du record sont
        protégés par un \texttt{std::mutex}, pris seulement quand un essaim bat le record.
        Les threads ne se ralentissent donc presque jamais.
  \item \textbf{Hasard reproductible par essaim.} Chaque thread a son propre générateur
        (\texttt{mt19937\_64}), initialisé à partir de \texttt{-{}-seed} et du numéro du thread.
  \item \textbf{Budget de temps.} Le temps est mesuré avec une horloge monotone
        (\texttt{steady\_clock}), insensible aux changements d'heure du système : 55
        secondes par défaut pour une limite d'une minute.
  \item \textbf{Rien n'est perdu.} Le meilleur planning est écrit dans le JSON au plus
        toutes les 2 secondes pendant la recherche, puis une dernière fois à la fin. Si le
        programme est interrompu, le dernier record est sur le disque.
  \item \textbf{Vérification.} \texttt{tools/validate.py} relit l'instance et le JSON, et
        contrôle les trois contraintes et le makespan, indépendamment du code C++.
\end{itemize}

\code{src/pso.cpp}{146}{163}{publier un nouveau record}
\begin{lignes}
  \item[\li{148}] test rapide sans verrou : si ce n'est pas mieux que le record, on s'arrête
        tout de suite. C'est le cas de presque toutes les évaluations.
  \item[\li{149--150}] sinon, on verrouille et on revérifie : entre-temps, un autre thread a
        peut-être trouvé mieux.
  \item[\li{151--153}] on construit le planning complet et on met à jour le record.
  \item[\li{154--156}] la ligne \texttt{[progression]} affichée dans le terminal.
  \item[\li{157--158}] sauvegarde du JSON et de la courbe de convergence (\texttt{-{}-log}).
  \item[\li{159--161}] avec \texttt{-{}-target}, on prévient tous les threads qu'ils peuvent s'arrêter.
\end{lignes}

\code{src/pso.cpp}{119}{125}{l'échéance}
\begin{lignes}
  \item[\li{119}] l'heure de départ, lue sur l'horloge monotone \texttt{steady\_clock}.
  \item[\li{122--125}] l'échéance = départ + \texttt{-{}-time} secondes ; avec
        \texttt{-{}-time 0}, il n'y a pas d'échéance. Chaque essaim teste
        \texttt{Clock::now() < deadline} à chaque tour (ligne 282).
\end{lignes}

\code{src/pso.cpp}{101}{115}{sauvegarder sans ralentir}
\begin{lignes}
  \item[\li{107--109}] si la dernière écriture date de moins de 2 secondes, on n'écrit pas
        (sauf \texttt{force}, utilisé à la toute fin).
  \item[\li{110--113}] le fichier s'appelle \texttt{<nom de l'instance>.json}, dans le dossier
        \texttt{-{}-out}, qui est créé s'il n'existe pas.
\end{lignes}

\code{src/main.cpp}{82}{92}{le programme principal}
\begin{lignes}
  \item[\li{82}] lecture de l'instance (section 3).
  \item[\li{84}] tout le travail : les essaims tournent jusqu'à l'échéance.
  \item[\li{87--88}] la ligne finale \texttt{makespan=\dots{} temps\_secondes=\dots}.
  \item[\li{89--91}] toute erreur (fichier introuvable, instance mal formée, option inconnue)
        est affichée proprement et le programme renvoie le code 1.
\end{lignes}

\paragraph{Le format de sortie.} Pour l'atelier $3\times3$ avec l'ordre A :
\begin{lstlisting}[style=json]
{
  "makespan": 11,
  "jobs_schedule": {
    "Job_0": [{"machine": 0, "start": 0, "end": 3},
              {"machine": 1, "start": 4, "end": 6},
              {"machine": 2, "start": 7, "end": 11}],
    "Job_1": [{"machine": 1, "start": 0, "end": 4},
              {"machine": 2, "start": 4, "end": 7},
              {"machine": 0, "start": 8, "end": 10}],
    "Job_2": [{"machine": 2, "start": 0, "end": 2},
              {"machine": 0, "start": 3, "end": 8},
              {"machine": 1, "start": 8, "end": 11}]
  }
}
\end{lstlisting}

\code{src/schedule.cpp}{143}{160}{écrire ce JSON}
\begin{lignes}
  \item[\li{143--144}] l'en-tête avec le makespan.
  \item[\li{145--146}] une entrée \texttt{"Job\_0"}, \texttt{"Job\_1"}\dots{} par job.
  \item[\li{151--156}] chaque étape devient \texttt{\{"machine", "start", "end"\}}, séparées par des virgules.
  \item[\li{158}] pas de virgule après le dernier job, sinon le JSON serait invalide.
\end{lignes}

% =====================================================================
\section{Résultats}
\label{sec:resultats}

Les mesures portent sur 7 instances de la bibliothèque JSPLIB, de $10\times10$ à
$30\times20$, dont l'optimum (ou la meilleure solution connue) est publié. Le critère est
l'\textbf{écart} au meilleur connu :
\[
\text{écart} = \frac{\text{makespan obtenu} - \text{meilleur connu}}{\text{meilleur connu}}.
\]

\subsection{Convergence sur l'instance des slides}

\begin{figure}[H]
\centering
\begin{tikzpicture}
\begin{axis}[width=0.92\textwidth,height=6.2cm,
  xmode=log,xmin=0.001,xmax=60,
  xlabel={temps (secondes, échelle logarithmique)},ylabel={meilleur makespan},
  ymin=1200,ymax=1700,
  grid=major,grid style={grille},
  legend style={font=\small,at={(0.98,0.96)},anchor=north east,draw=none,fill=white,fill opacity=0.85},
  every axis plot/.append style={const plot,very thick}]
  \addplot[gris,dashed] table[col sep=semicolon,x expr=\thisrow{temps_ms}/1000+0.001,y=meilleur_makespan] {donnees/conv_base.csv};
  \addlegendentry{PSO de base (semi-actif, global, 1 essaim)}
  \addplot[accent] table[col sep=semicolon,x expr=\thisrow{temps_ms}/1000+0.001,y=meilleur_makespan] {donnees/conv_final.csv};
  \addlegendentry{PSO final (GT, mixte, multi-essaims)}
  \addplot[jdeux,thick,dotted,no markers,const plot=false] coordinates {(0.001,1234) (60,1234)};
  \addlegendentry{optimum connu (1234)}
\end{axis}
\end{tikzpicture}
\caption{Évolution du meilleur makespan sur abz5 (l'instance \texttt{10\_10} des slides),
en 55 secondes. Chaque marche est un nouveau record. Le PSO de base trouve son dernier
record (1283, soit 4,0~\% au-dessus de l'optimum) après 5,4~s puis n'avance plus. Le PSO
final atteint 1239 (0,4~\%) dès 0,25~s : sur une instance de cette taille, ses 12 essaims
trouvent très vite un planning proche de l'optimum, puis plafonnent eux aussi.}
\label{fig:convergence}
\end{figure}

\subsection{Ce qu'apporte chaque mécanisme}

Chaque ligne ajoute un mécanisme à la précédente. Les mesures sont faites en 10 secondes
sur 2 à 3 graines aléatoires, ce qui permet de comparer rapidement. La dernière ligne est
mesurée sur le vrai budget de 55 secondes.

\begin{table}[h]
\centering
\caption{Écart moyen au meilleur connu sur les 7 instances.}
\label{tab:ablation}
\begin{tabular}{lr}
\toprule
Configuration & Écart moyen\\
\midrule
Décodeur semi-actif, topologie globale (version de départ) & 22,3 \%\\
Ancien décodeur \og GT \fg{} approximatif & 10,4 \%\\
Nouveau décodeur Giffler-Thompson paramétré & 6,8 \%\\
\quad + topologie en anneau & 6,4 \%\\
\quad + redémarrage après 200 tours sans progrès & 5,9 \%\\
\quad + voisinage de 2 et $\delta=0{,}2$ & 5,1 \%\\
\quad + 12 essaims indépendants (un par cœur) & 4,7 \%\\
\quad + migration du record toutes les 50 itérations & 4,2 \%\\
\quad + essaims mixtes anneau / CLPSO & 3,8 \%\\
\midrule
\textbf{Configuration finale, 55 secondes} & \textbf{3,3 \%}\\
\bottomrule
\end{tabular}
\end{table}

\subsection{Résultats par instance}

\begin{table}[h]
\centering
\caption{Version de départ (10 s, 1 essaim) contre version finale (55 s, 12 essaims).}
\label{tab:instances}
\begin{tabular}{llrrrr}
\toprule
Instance & Taille & Meilleur connu & Départ & Final & Makespan final\\
\midrule
abz5 (\texttt{10\_10}) & $10\times10$ & 1234 & 3,8 \% & 0,4 \% & 1239\\
ft10 & $10\times10$ & 930  & 7,5 \%  & 2,0 \%  & 949\\
la21 & $15\times10$ & 1046 & 9,4 \%  & 0,9 \%  & 1055\\
la26 & $20\times10$ & 1218 & 8,4 \%  & \textbf{0,0 \%} & \textbf{1218}\\
ta01 & $15\times15$ & 1231 & 10,4 \% & 4,1 \%  & 1282\\
ta21 & $20\times20$ & 1644 & 14,7 \% & 5,2 \%  & 1729\\
ta41 & $30\times20$ & 2018 & 18,5 \% & 10,7 \% & 2234\\
\midrule
\textbf{Moyenne} & & & \textbf{10,4 \%} & \textbf{3,3 \%} &\\
\bottomrule
\end{tabular}
\end{table}

\begin{table}[h]
\centering
\caption{Anneau contre CLPSO seuls (55 s, 1 essaim, moyenne de 3 graines) : des profils
complémentaires, d'où le choix d'essaims mixtes.}
\label{tab:ringcl}
\begin{tabular}{lrrrrrrr}
\toprule
 & abz5 & ft10 & la21 & la26 & ta01 & ta21 & ta41\\
\midrule
Anneau + redémarrage & 0,41 & 2,15 & 2,96 & 0,93 & 4,93 & \textbf{6,97} & \textbf{12,09}\\
CLPSO                & 0,49 & 2,15 & \textbf{1,15} & \textbf{0,00} & \textbf{4,01} & 7,89 & 14,87\\
\bottomrule
\end{tabular}
\end{table}

\subsection{Ce que nous avons appris}

\begin{itemize}
  \item \textbf{Le décodeur compte autant que l'essaim.} Changer de décodeur a fait plus
        que tous les réglages d'essaim réunis : c'est lui qui définit l'espace où les
        particules cherchent.
  \item \textbf{La stagnation se combat par la diversité, pas par le temps.} Doubler le
        temps de calcul n'a presque rien changé (5,17~\% $\to$ 5,09~\%), alors que les
        topologies, les redémarrages et les essaims multiples ont fait gagner plusieurs
        points.
  \item \textbf{Il faut mesurer avant de garder une idée.} La relance autour du record
        semblait prometteuse ; elle faisait moins bien qu'une relance aléatoire. Les petits
        réglages (nombre de particules, inertie, \texttt{vmax}) donnaient des écarts plus
        petits que le bruit de mesure ($\pm0{,}3$~\%).
  \item \textbf{Limite.} Les très grandes instances restent difficiles (ta41 à 10,7~\%) :
        avec 600 clés par particule, un PSO pur peine à explorer un espace aussi vaste en
        une minute.
\end{itemize}

% =====================================================================
\section{Utilisation}

\paragraph{Compiler} (une fois, en mode optimisé) :
\begin{lstlisting}[style=shell]
cmake -S . -B build -DCMAKE_BUILD_TYPE=Release
cmake --build build
ctest --test-dir build --output-on-failure
\end{lstlisting}

\paragraph{Lancer une instance} avec les réglages par défaut, qui sont les meilleurs trouvés :
\begin{lstlisting}[style=shell]
./build/jsp_pso instances/abz5 --out resultats
python3 tools/validate.py instances/abz5 resultats/abz5.json
\end{lstlisting}

\begin{table}[h]
\centering
\caption{Les options de la ligne de commande.}
\small
\begin{tabular}{>{\ttfamily}l l p{8.3cm}}
\toprule
\textrm{Option} & Défaut & Rôle\\
\midrule
-{}-time S       & 55     & budget en secondes (0 = sans limite)\\
-{}-threads N    & 0      & nombre d'essaims, 0 = un par cœur\\
-{}-topology     & mixed  & \texttt{ring}, \texttt{global}, \texttt{cl} (CLPSO) ou \texttt{mixed}\\
-{}-neighbors K  & 2      & voisins de chaque côté dans l'anneau\\
-{}-restart N    & 200    & tours sans progrès avant relance (0 = jamais)\\
-{}-migrate N    & 50     & période de migration du record (0 = jamais)\\
-{}-particles N  & 30     & particules par essaim\\
-{}-vmax V       & 0,2    & vitesse maximale par dimension\\
-{}-decoder      & gt     & \texttt{gt} (Giffler-Thompson) ou \texttt{semi} (semi-actif)\\
-{}-delta D      & 0,2    & fenêtre du décodeur \texttt{gt}\\
-{}-seed N       & 0      & graine du hasard\\
-{}-out DIR      & .      & dossier du JSON de sortie\\
-{}-target M     &        & s'arrêter dès qu'un makespan $\le M$ est trouvé\\
-{}-log FICHIER  &        & écrire la courbe de convergence (CSV)\\
\bottomrule
\end{tabular}
\end{table}

\paragraph{Tests automatiques.} \texttt{ctest} vérifie la lecture de l'exemple du cours et
de huit fichiers volontairement défectueux, le décodage de l'ordre A (makespan 11),
l'écriture du JSON, et que le PSO retrouve l'optimum 11 sur l'atelier $3\times3$. Une
compilation avec \emph{sanitizers} (\texttt{-DJSP\_SANITIZE=ON}) détecte en plus les erreurs
mémoire et les comportements indéfinis.

\end{document}
